\documentclass[11pt]{article}

\usepackage[french, english]{babel}

\usepackage[T1]{fontenc}

\usepackage{amsfonts}
\usepackage{amsmath}
\usepackage{amssymb}

\usepackage{graphicx}
\usepackage{multirow}
\usepackage{multicol}
\usepackage{xcolor}
\usepackage[shortlabels]{enumitem}

\usepackage[margin=2cm]{geometry}

\usepackage{listings}

\usepackage{tikz}
\usetikzlibrary{shapes.geometric, arrows.meta, positioning}

\tikzstyle{block} = [rectangle, rounded corners, minimum width=4cm, minimum height=1cm,
text centered, draw=black, fill=blue!10]
\tikzstyle{process} = [rectangle, rounded corners, minimum width=5cm, minimum height=1cm,
text centered, draw=black, fill=green!10]
\tikzstyle{io} = [trapezium, trapezium left angle=70, trapezium right angle=110,
minimum width=4cm, minimum height=1cm, text centered, draw=black, fill=orange!15]
\tikzstyle{decision} = [diamond, aspect=2, draw=black, fill=yellow!15, text centered]
\tikzstyle{arrow} = [thick, ->, >=Stealth]

\renewcommand\lstlistingname{Code Python}

\lstset{frame=tb,
  language=Python,
  numbers=left,
  stepnumber=1,
  aboveskip=3mm,
  belowskip=3mm,
  showstringspaces=false,
  columns=flexible,
  basicstyle={\small\ttfamily},
%   numbers=none,
  numberstyle=\tiny\color{gray},
  keywordstyle=\color{blue},
  commentstyle=\color{paired2b},
  stringstyle=\color{paired3b},
  breaklines=true,
  breakatwhitespace=true,
  tabsize=1,
  frame=single,
%   showtabs=false,
}

\begin{document}
\begin{multicols}{2}
  \section{\underline{Calculabilité}}
  \subsection{Rappels}
  \textbf{Automate} : modèle mathématique de calcul.\\
  \textbf{Automate fini déterministe (AFD)} : automate avec un seul état possible à chaque instant.\\
  \textbf{Automate fini non déterministe (AFN)} : automate avec plusieurs états possibles à chaque instant.\\
  \textbf{Alphabet} : ensemble fini t.q. $\Sigma=\{0,1\}$\\
  \textbf{Mot} : suite finie de symboles d'un alphabet. Ex : $w=0110 \in \Sigma^4$\\
  \textbf{Langage} : ensemble de mots sur un alphabet. Ex : $L=\{w | w \text{ commence par } 0\}$\\
  \textbf{Taille} : nombre de symboles d'un mot. Ex : $|w|=4$\\\textbf{Concatenation} : opération binaire entre deux mots. Ex : $w_1=01, w_2=10 \Rightarrow w_1w_2=0110$\\
  \textbf{Concatenation} : opération binaire entre deux mots. Ex : $w_1=01, w_2=10 \Rightarrow w_1w_2=0110$\\
  \textbf{Ensemble des langages} : ensemble des parties de $\Sigma^*$. Ex : $\mathcal{L}=\mathcal{P}(\Sigma^*)$\\
  \textbf{Renversé} : mot lu à l'envers. Ex : $w=0110 \Rightarrow w^R=0110$\\

  \subsection{Machines de Turing}
  \textbf{Définition} : une machine de Turing est un 7-uplet $(Q, \Sigma, \Gamma, \delta, q_0, q_{accept}, q_{reject})$ où :\\
  - $Q$ est un ensemble fini d'états\\
  - $\Sigma$ est un alphabet d'entrée ne contenant pas le symbole blanc $\sqcup$\\
  - $\Gamma$ est un alphabet de bande, avec $\Sigma \subset \Gamma$\\
  - $\delta : Q \times \Gamma \rightarrow Q \times \Gamma \times \{L, R\}$ est la fonction de transition\\
  - $q_0 \in Q$ est l'état initial\\
  - $q_{accept} \in Q$ est l'état d'acceptation\\
  - $q_{reject} \in Q$ est l'état de rejet, avec $q_{reject} \neq q_{accept}$\\
\end{multicols}
\end{document}